\makeabstract
{ % #1: Title
Electrophysiological Diversity of Mouse Vestibular Ganglion Glia
}
{ % #2: Authorship
Hannah Feng\textsuperscript{1},
Radha Kalluri\textsuperscript{2}
}
{ % #3: Affiliations (use \\ to start a newline)
\textsuperscript{1} Amherst College, Amherst, MA\\
\textsuperscript{2} Department of Otolaryngology, University of Southern California Keck School of Medicine, Los Angeles, CA
}
{ % #4: Purpose
In the peripheral nervous system, glial cells are responsible for the formation and maintenance of myelin, which is critical for signal propagation and neuronal survival. Beyond their role in myelination, little is known about the function of vestibular glia. Recent work in mouse spiral ganglion glia identified both inwardly and outwardly rectifying currents around the onset of hearing, suggesting a diversification of cochlear glial properties. Elsewhere in the nervous system, glia have also been shown to clear neurotransmitters, buffer potassium, and provide energy to neurons. As the vestibular ganglion and the spiral ganglion deliver different information from the ear to the brain, my work will investigate whether glia in the vestibular ganglion exhibit a similar electrophysiological diversity to those of the spiral ganglion.
}
{ % #5: Methods
Vestibular ganglia were excised from four C57BL/6 mice postnatal day 9 to 16. The ganglia were enzymatically treated, dissociated, and cultured for 1-2 nights to isolate single somata for patch clamp recordings. Recordings were made in the rupture-patch configuration to characterize whole-cell currents. To assess their voltage-activated membrane currents, glia were subjected to series of 200 ms voltage steps from -203 to +57 mV, in 10 mV increments, from a holding potential of -73 mV. Additionally, glial fibrillary acidic protein (GFAP) and K\textsubscript{ir}4.1, an inwardly rectifying potassium channel, were detected using immunofluorescence in the vestibular ganglion of a postnatal day 29 mouse.
}
{ % #6: Results
Vestibular ganglion glia comprised distinct inwardly and outwardly rectifying cell populations, with inwardly rectifying currents being larger in amplitude. Of 26 successful glia with giga-ohm seals, 20/26 showed inwardly rectifying currents, and 6/26 showed outwardly rectifying currents. Glia with inwardly rectifying currents also had a more hyperpolarized resting membrane potential at -72.8 mV, close to the reversal potential of potassium. K\textsubscript{ir}4.1 immunodetection showed strong fluorescence in the satellite glial cell layer but weak fluorescence in the Schwann cell layer.
}
{ % #7: Conclusion
In both the spiral ganglion and the vestibular ganglion, glia exhibit diverse membrane properties, indicating a divergence in function. For satellite glial cells, hypothesized to be the population of glia with inwardly rectifying currents, potassium clearance may be a critical role during the neuron's repolarization step. For Schwann cells, hypothesized to be the cells with outwardly rectifying currents, outward K\textsubscript{v} currents have been reported to influence cell proliferation and differentiation. Mammalian Schwann cells have been shown to express K\textsubscript{ir} only around the onset of myelination, which may indicate other mechanisms of potassium clearance following maturation.
}
 \newpage